\section{A separation-dependent marked pair statistic}\label{sec:pairs}

The reduction in Proposition~\ref{prop:quadratic} motivates retaining
both transverse marks separately. The following test is nonnegative,
vanishes on real configurations, and includes a nonconstant separation
kernel. Its approximation follows from the same one-variable logarithm
estimate, without introducing a probability law.

\subsection{A fixed test and its exact response}

Define the piecewise polynomial cutoffs
\begin{equation}\label{eq:p-cutoff}
 p(t)=
 \begin{cases}
 1,&t\le0,\\
 1-10t^3+15t^4-6t^5,&0<t<1,\\
 0,&t\ge1,
 \end{cases}
 \qquad w(t)=p(2|t|-1),\qquad v(t)=t_+^2p(t-1).
\end{equation}
Thus $w=1$ on $[-1/2,1/2]$, $w=0$ off $[-1,1]$,
and $v(t)=t^2$ for $0\le t\le1$, with $v=0$ off $[0,2]$.
The function $v$ is $C^{1,1}$; smoothness of higher order at zero
is neither asserted nor needed.

For $X\ge10$ and $0<h,b\le1$, let
\begin{align}
 a(z)&=w((x-X)/h)v(y/b),&
 q(z)&=w((x-X)/h)v(|y|/b),\label{eq:pair-marks}\\
 c_h(u)&=\frac{h^2}{u^2+h^2},&
 K(z,z')&=a(z)a(z')c_h(x-x').\notag
\end{align}
For any conjugation-symmetric holomorphic field defined on a
neighborhood of the support of $q$, set
\begin{equation}\label{eq:marked-statistic}
 S_{X,h,b}(H)=\sum_{j\ne k}q(z_j)q(z_k)c_h(x_j-x_k).
\end{equation}
All multiplicity copies are separate labels. The center is fixed;
there is no implicit averaging over $X$.

\begin{proposition}[Marked-pair signal]\label{prop:pair-signal}
Every term in \eqref{eq:marked-statistic} is nonnegative.
A simple conjugate pair $x_*\pm\ii d$ with
$|x_*-X|\le h/2$ and $0<d\le b$ forces
\begin{equation}\label{eq:marked-signal}
 S_{X,h,b}(H)\ge2(d/b)^4 .
\end{equation}
For $m$ copies at each conjugate location, their population contributes
$(4m^2-2m)(d/b)^4$. Every all-real configuration gives zero,
regardless of multiplicity.
\end{proposition}
\begin{proof}
On the plateau each of the $2m$ labels has weight $(d/b)^2$.
Every difference of their horizontal coordinates is zero, so
$c_h=1$. There are $2m(2m-1)$ ordered distinct-label pairs.
Other summands cannot reduce their contribution. Real labels
have $q=0$.
\end{proof}

A useful declared scale is
\begin{equation}\label{eq:matched-scale}
 h=b=h_X:=\min(1,4\pi/\log X).
\end{equation}
This uses the leading mean-spacing scale in the $x=2\gamma$
coordinate; it is a definition, not an exact spacing assertion.
The band $b/2\le d\le b$ has signal at least $1/8$.
A single such band does not cover arbitrarily small transverse
displacements or all displacements up to the canonical strip edge.

\subsection{Conjugate folding and the self-index correction}

The support of $a$ is contained in $D(c,r)$ with the same buffer as
in Appendix~\ref{app:scalar}. For a nonzero holomorphic $V$ on
the disk, write $\mu_V=\nu_V$ there and define
\[
 T_V(K)=\iint K\dd\mu_V\dd\mu_V,\qquad
 D_V(a^2)=\int a^2\dd\mu_V .
\]
If $V$ has the same zeros on this disk as a conjugation-symmetric
full field, then
\begin{equation}\label{eq:folded-pair}
 S(V)=4T_V(K)-2D_V(a^2).
\end{equation}
Each upper nonreal label represents itself and its lower conjugate,
so the full tensor square has factor four. Removing the self label
of each of these two copies has factor two. At an upper point of
multiplicity $m$, the tensor has mass $m^2$ and the one-point
correction removes $m$, rather than $m^2$. In particular the two
conjugate partners are different full-plane labels even when their
upper representatives coincide.

Applying \eqref{eq:weak-onepoint} separately in the two independent
variables gives
\begin{align}
 T_V(K)&=\frac1{4\pi^2}
  \iint\log|V(z)|\log|V(z')|\,
       \Delta_z\Delta_{z'}K(z,z')\dd A\dd A',\label{eq:pair-log}\\
 D_V(a^2)&=\frac1{2\pi}
       \int\log|V|\,\Delta(a^2)\dd A .\notag
\end{align}
For this particular $K$ the mixed weak derivative is bounded:
each factor $a$ is differentiated at most twice. Product logarithms
are integrable by Fubini, including near the geometric diagonal,
because they belong to different variables. Mollification in the
support buffer proves the formula for the stated cutoffs.
Adding a harmonic term to either logarithm changes neither integral;
integrate first in its corresponding variable. Thus scalar
normalizations and zero-free multipliers are immaterial.

\subsection{Tensor transfer with only second derivatives}

\begin{theorem}[Deterministic pair transfer]\label{thm:pair-transfer}
Let $f,g$ satisfy the normalized disk and anchor hypotheses of
Theorem~\ref{thm:log}, with the disk translated to $c$, and suppose
\[
 \int_{D(c,r)}|\log|f|-\log|g||\dd A\le E .
\]
Let $K$ be a real symmetric $C_c^{1,1}$ test on $D(c,r)^2$.
Its sections and diagonal trace have bounded weak Laplacians.
Choose bounds
\[
 D_2\ge\max\{\norm{\Delta_zK}_\infty,
                  \norm{\Delta_{z'}K}_\infty\},\qquad
 D_{\rm diag}\ge\norm{\Delta[K(z,z)]}_\infty .
\]
For the folded functional
$\mathcal S_V(K)=4T_V(K)-2\int K(z,z)\dd\mu_V(z)$,
\begin{equation}\label{eq:pair-transfer}
 |\mathcal S_f(K)-\mathcal S_g(K)|
 \le\frac{E}{\pi}\bigl(4nD_2+D_{\rm diag}\bigr),
\end{equation}
where $n$ is \eqref{eq:local-constants}. For the test
\eqref{eq:pair-marks} and conjugation-symmetric full fields,
this is the difference of their statistics
\eqref{eq:marked-statistic}.
\end{theorem}
\begin{proof}
Each relevant population is at most $n$ by the Jensen part of
Theorem~\ref{thm:log}. For fixed $z'$ in the relevant projection
of the support, the weak one-variable identity gives
\[
 |(\mu_f-\mu_g)(K(\,\cdot\,,z'))|
 \le\frac{ED_2}{2\pi}.
\]
The exact telescoping identity
\[
 \mu_f\otimes\mu_f-\mu_g\otimes\mu_g
 =(\mu_f-\mu_g)\otimes\mu_f+
                         \mu_g\otimes(\mu_f-\mu_g)
\]
allows us to integrate this bound against positive zero measures.
The two masses cost at most $2n$, giving
$|T_f-T_g|\le nED_2/\pi$.
The diagonal trace is a one-variable test, so its error is at most
$ED_{\rm diag}/(2\pi)$. Multiplication by four and two proves
\eqref{eq:pair-transfer}.
\end{proof}

The population bound replaces a fourth-derivative cost. Indeed if
$D_4\ge\norm{\Delta_z\Delta_{z'}K}_\infty$ and
$M=\pi r^2A$, \eqref{eq:log-mass} and
\[
 L_f(z)L_f(z')-L_g(z)L_g(z')
 =(L_f-L_g)(z)L_f(z')+L_g(z)(L_f-L_g)(z')
\]
give the valid alternative
\begin{equation}\label{eq:pair-fourth}
 |\mathcal S_f-\mathcal S_g|
 \le E\left(\frac{2MD_4}{\pi^2}
                      +\frac{D_{\rm diag}}{\pi}\right).
\end{equation}
Theorem~\ref{thm:pair-transfer} instead uses only second derivatives.
This compares two explicit allowances, not optimal approximation
rates. Strong deterministic accuracy on the same disk suffices
for this transfer; an independently favorable pair value remains
a different question.

\subsection{Explicit constants and the canonical theorem}

Differentiating \eqref{eq:p-cutoff},
$p'=-30t^2(1-t)^2$ and
$p''=-60t(1-t)(1-2t)$ on $(0,1)$.
Their maxima satisfy
$\norm{p'}_\infty\le15/8$ and
$\norm{p''}_\infty\le10/\sqrt3<6$.
Consequently
\[
 \norm w_\infty\le1,\quad \norm{w'}_\infty\le15/4,\quad
 \norm{w''}_\infty\le24,\qquad
 \norm v_\infty\le4,\quad \norm{v'}_\infty\le23/2,\quad
 \norm{v''}_\infty\le41.
\]
For example on $1<t<2$,
$v'=2tp+t^2p'$ and
$v''=2p+4tp'+t^2p''$, which give the last two bounds.
Thus
\begin{equation}\label{eq:a-bounds}
 |a|\le A_0=4,\quad |a_x|\le A_x=15/h,\quad
 |a_y|\le A_y=23/(2b),\quad
 |\Delta a|\le D_a=96h^{-2}+41b^{-2}.
\end{equation}
The partial fraction formula
$c_h(u)=\frac{h}{2\ii}[(u-\ii h)^{-1}-(u+\ii h)^{-1}]$
implies $|c_h^{(j)}(u)|\le j!h^{-j}$ on the real axis.
Since
\[
 \Delta_zK=(\Delta a)a'c_h+2a_xa'c_h'+aa'c_h'',
 \qquad K(z,z)=a(z)^2,
\]
where $a'=a(z')$ in this display only, admissible constants are
\begin{align}
 D_2&=A_0D_a+2A_xA_0/h+2A_0^2/h^2
       =536h^{-2}+164b^{-2},\label{eq:D2}\\
 D_{\rm diag}
   &=2A_0D_a+2A_x^2+2A_y^2
       =1218h^{-2}+\frac{1185}{2}b^{-2}.\label{eq:Ddiag}
\end{align}
For \eqref{eq:pair-fourth}, expansion of the mixed Laplacian permits
the choice
\[
 D_4= D_a^2+\frac{4A_xD_a}{h}
 +\frac{4A_0D_a+8A_x^2}{h^2}
 +\frac{24A_xA_0}{h^3}+\frac{24A_0^2}{h^4}
 =D_a^2+76D_a/h^2+3624/h^4.
\]
No fourth derivative of $v$ has been used.

\begin{theorem}[Complete canonical marked-pair acquisition]\label{thm:canonical-pair}
Let $X\ge10$, $0<h,b\le1$, and define the test by
\eqref{eq:p-cutoff}--\eqref{eq:marked-statistic}.
Choose finite $J,U$ so that the complete error
\eqref{eq:complete-delta} satisfies \eqref{eq:delta-requirement}.
Then
\begin{align}
 |S_{X,h,b}(H_0)-S_{X,h,b}(\widehat H_{J,U})|
 &\le E_{\rm pair},\label{eq:canonical-pair}\\
 E_{\rm pair}
 &=\frac{C_X(\delta/M_X)^{\alpha_X}}{\pi}
 \left[
 4n_X(536h^{-2}+164b^{-2})
 +1218h^{-2}+\frac{1185}{2}b^{-2}
 \right].\notag
\end{align}
In particular, when $b=h$,
\begin{equation}\label{eq:matched-error}
 \boxed{\displaystyle
 E_{\rm pair}=\frac{2800n_X+3621/2}{\pi h^2}\,
                  C_X(\delta/M_X)^{2/(n_X+2)} .}
\end{equation}
The constants and both omitted source tails are explicitly
specified in Section~\ref{sec:canonical}.
\end{theorem}
\begin{proof}
Normalize $F$ and $\widehat F$ as in \eqref{eq:canonical-constants}.
Equation~\eqref{eq:EX} gives $E=E_X$, and their zero populations are
at most $n_X$. Apply Theorem~\ref{thm:pair-transfer} with
\eqref{eq:D2}--\eqref{eq:Ddiag}. Both original cosine fields are
conjugation-symmetric; $D$ is zero-free on the disk. Thus the
folded functionals equal the stated full-plane statistics.
The approximant's zeros in $b<|y|<2b$ retain their actual weights.
For $b=h$, $4n_X(536+164)=2800n_X$ and
$1218+1185/2=3621/2$.
\end{proof}

At the scale \eqref{eq:matched-scale}, the prefactor multiplying
$(\delta/M_X)^{\alpha_X}$ is $O((\log X)^5)$.
The direct log-product allowance \eqref{eq:pair-fourth} instead
has prefactor $O((\log X)^7)$. These are upper-bound costs, not
actual error asymptotics or runtime estimates; both retain the
deteriorating exponent and the enlarged-disk source cost.

Write $K_{\rm pair}=(4n_XD_2+D_{\rm diag})/\pi$.
For any $e>0$, use \eqref{eq:precision-choice} with
$K_\varphi$ replaced by $K_{\rm pair}$ and then
\eqref{eq:tail-allocation}; this gives $E_{\rm pair}\le e/2$.
If a later calculation obtains $\widehat q$ with
$|\widehat q-S(\widehat H_{J,U})|\le E_{\rm eval}$, the precise
condition excluding a plateau pair with $d_0\le d\le b$ is
\begin{equation}\label{eq:pair-consumer}
 \widehat q+E_{\rm eval}+E_{\rm pair}<2(d_0/b)^4 .
\end{equation}
At present the favorable value on the left has not been acquired.

\subsection{Two scale checks}

A center average over an interval of length $L$ containing the
whole pair-window translate dilutes its mutual contribution to
\[
 2(d/b)^4\frac{h}{L}\int_\R w(t)^2\dd t
 =2(d/b)^4\frac{h}{L}\,\frac{643}{462}.
\]
The constant is elementary:
$\int_\R w^2=1+\int_0^1p(t)^2\dd t=1+181/462$.
Thus the fixed-center $1/8$ band margin cannot be carried
unchanged into an averaged theorem.

The scale dependence of the transfer also cannot be discarded.
For $0<d\le1$ consider the local polynomial controls
\begin{equation}\label{eq:pair-control}
 f_\pm(z)=\frac{(z-X)^2\pm d^2}{74}.
\end{equation}
On $\overline D(c,R)$, $|z-X|\le17/2$, so
$\norm{f_\pm}_\infty<1$,
$|f_\pm(c)|\ge15/74$ and
$\norm{f_+-f_-}_\infty=d^2/37$.
Yet $f_+$ has the conjugate pair $X\pm\ii d$ and $f_-$ has
only real roots. Taking $h=b=2d\le1$ gives
$S(f_+)=1/8$ and $S(f_-)=0$ while the field difference tends
to zero. A modulus of continuity independent of the shrinking
test scale is therefore impossible for this generic class.
The controls are local polynomials, not canonical theta fields;
they neither refute Theorem~\ref{thm:canonical-pair} nor decide
a source-specific favorable estimate.
